% Enterprise LLM TechDoc LaTeX Template (ctexart)
\documentclass[UTF8,a4paper,11pt,fontset=fandol]{ctexart}

% === Core packages ===
\usepackage{geometry}
\usepackage{booktabs}
\usepackage{array}
\usepackage{tabularx}
\usepackage{graphicx}
\usepackage{amsmath,amssymb}
\usepackage{hyperref}
\usepackage{xcolor}
\usepackage{caption}
\usepackage{float}
\usepackage{enumitem}
\usepackage{fancyhdr}
\usepackage{tcolorbox}
\usepackage{tikz}
\usetikzlibrary{arrows.meta,positioning,shapes.misc,shapes.geometric}

\geometry{a4paper, top=2.3cm, bottom=2.3cm, left=2.2cm, right=2.2cm}
\setlength{\parskip}{0.35em}
\setlength{\parindent}{2em}
\renewcommand{\arraystretch}{1.25}

% === Colors ===
\definecolor{Primary}{HTML}{1A1A2E}
\definecolor{Accent}{HTML}{1F6FEB}
\definecolor{Muted}{HTML}{6B7280}
\definecolor{BoxBg}{HTML}{F8FAFC}

% === Hyperref ===
\hypersetup{
  colorlinks=true,
  linkcolor=Accent,
  urlcolor=Accent,
  citecolor=Accent,
  pdftitle={大模型能力与现有产品结合技术方案（模板样例）},
  pdfauthor={解决方案团队（示例）}
}

% === Header / Footer ===
\pagestyle{fancy}
\fancyhf{}
\lhead{\textcolor{Muted}{某某客户（示例）}}
\rhead{\textcolor{Muted}{V0.1}}
\cfoot{\textcolor{Muted}{\thepage}}
\renewcommand{\headrulewidth}{0.3pt}
\renewcommand{\footrulewidth}{0pt}

% === Boxes ===
\newtcolorbox{infobox}[1]{
  colback=BoxBg,
  colframe=Accent!45,
  title=\textbf{#1},
  boxrule=0.6pt,
  arc=2pt,
  left=6pt,right=6pt,top=6pt,bottom=6pt
}

% === Helpers ===
\newcolumntype{L}[1]{>{\raggedright\arraybackslash}p{#1}}
\newcolumntype{Y}{>{\raggedright\arraybackslash}X}

\title{\textbf{大模型能力与现有产品结合技术方案（模板样例）}\\[6pt]\large\textcolor{Muted}{客户技术文档（用于校验排版与结构）}}
\author{解决方案团队（示例）}
\date{2026-04-13}

\begin{document}
\maketitle

\begin{infobox}{文档控制信息}
\begin{tabularx}{\linewidth}{L{3.2cm}Y L{3.0cm}Y}
\toprule
客户名称 & 某某客户（示例） & 行业 & 零售/电商（示例） \\
版本 & V0.1 & 日期 & 2026-04-13 \\
作者 & 解决方案团队（示例） & 审阅 & 交付/安全（示例） \\
保密级别 & 内部/客户专用（示例） &  &  \\
\bottomrule
\end{tabularx}
\end{infobox}

\vspace{0.5em}
\tableofcontents
\newpage

\section{概述}
\subsection{目标与范围}
本文档用于描述\textbf{目标大模型服务（以合同与实际接入平台为准）}能力与客户现有产品/系统的结合方式，输出可交付的技术方案结构（排版与内容为模板样例）。对任何具体能力、指标与 SLA 的承诺以合同与实际开通能力为准。

\subsection{术语与约束}
\begin{itemize}[leftmargin=2.2em]
  \item 大模型能力范围：仅覆盖已开通/已采购能力；不确定项标记“待确认”。
  \item 数据合规：遵循数据分级分类、最小权限、脱敏与审计要求。
  \item 输出边界：本模板强调“能力—系统—场景—KPI”的可追溯关系，避免泛化营销表述。
\end{itemize}

\section{客户现有产品与集成切入点}
\begin{infobox}{现有产品清单（示例）}
\begin{tabularx}{\linewidth}{L{3.5cm} L{2.7cm} Y}
\toprule
产品/系统 & 责任部门 & 现状与说明 \\
\midrule
客服工单系统 & 客服中心 & 工单流转、知识库查询、质检 \\
企业知识库/文档平台 & 技术/运营 & 制度、FAQ、产品说明 \\
CRM/会员系统 & 销售/运营 & 客户画像、标签、活动触达 \\
APP/小程序 & 产品 & 用户自助服务入口 \\
\bottomrule
\end{tabularx}
\end{infobox}

\section{大模型能力描述（按可交付能力组织）}
\begin{infobox}{能力清单（示例，具体以实际开通为准）}
\begin{tabularx}{\linewidth}{L{5.2cm} Y}
\toprule
能力 & 说明与典型用途 \\
\midrule
对话式问答与多轮交互 & 用于客服、自助服务、员工助手 \\
长文理解与摘要/要点提取 & 用于工单归因、会议纪要、制度解读 \\
知识检索增强生成（RAG） & 基于企业知识库的可追溯回答 \\
信息抽取与结构化 & 从工单/聊天中抽取字段、意图、情绪 \\
内容生成与改写 & 营销文案、话术、公告模板（需合规审核） \\
工具调用/流程编排（Agent） & 调用工单、CRM、知识库检索、订单查询等 \\
\bottomrule
\end{tabularx}
\end{infobox}

\section{能力与现有产品结合分析（映射与收益）}
\subsection{映射方法}
建议采用“\textbf{场景}→\textbf{能力}→\textbf{集成点}→\textbf{收益}→\textbf{KPI}”的结构化映射，便于后续验收与运营复盘。

\subsection{能力—系统—场景映射表（模板）}
\begin{infobox}{映射表（示例）}
\begin{tabularx}{\linewidth}{L{3.3cm} L{3.2cm} L{3.4cm} Y}
\toprule
场景 & 大模型能力 & 集成系统/接口 & 预期收益与KPI \\
\midrule
场景A：智能客服与工单自动化 & （示例）见能力清单 & 工单系统API、知识库检索、订单/物流查询接口、客服坐席工作台 & 提升自助解决率，降低人工接入量、提升工单分类一致性与质检覆盖率；KPI：自助解决率、平均响应时长、人工转接率、工单一次分类准确率 \\
场景B：员工知识助手（制度/产品/流程） & （示例）见能力清单 & 文档平台/知识库、权限系统（SSO/LDAP）、企业IM/门户 & 提升信息获取效率、降低重复咨询与培训成本；KPI：平均检索耗时、员工满意度、重复咨询量 \\
\bottomrule
\end{tabularx}
\end{infobox}

\section{业务场景方案与优势（按场景输出）}
\subsection{场景A：智能客服与工单自动化}
\begin{infobox}{痛点}
\begin{itemize}[leftmargin=2.2em]
  \item 人工响应压力大，重复问题占比高
  \item 工单分类不一致，质检与复盘成本高
\end{itemize}
\end{infobox}
\begin{infobox}{方案要点}
以对话式大模型 + RAG 为核心，通过工具调用对接工单系统与订单查询，实现自助解答、自动建单、自动分类与质检辅助。
\end{infobox}
\begin{infobox}{集成与指标}
\begin{itemize}[leftmargin=2.2em]
  \item 集成点：工单系统API、知识库检索、订单/物流查询接口、客服坐席工作台
  \item 预期收益：提升自助解决率，降低人工接入量、提升工单分类一致性与质检覆盖率
  \item KPI：自助解决率、平均响应时长、人工转接率、工单一次分类准确率
\end{itemize}
\end{infobox}

\subsection{场景B：员工知识助手（制度/产品/流程）}
\begin{infobox}{痛点}
\begin{itemize}[leftmargin=2.2em]
  \item 制度分散、更新频繁，新员工学习成本高
  \item 跨部门协作时信息检索耗时
\end{itemize}
\end{infobox}
\begin{infobox}{方案要点}
将企业知识库接入 RAG，提供可引用来源的问答，并结合摘要与要点提取输出「结论 + 依据」。
\end{infobox}
\begin{infobox}{集成与指标}
\begin{itemize}[leftmargin=2.2em]
  \item 集成点：文档平台/知识库、权限系统（SSO/LDAP）、企业IM/门户
  \item 预期收益：提升信息获取效率、降低重复咨询与培训成本
  \item KPI：平均检索耗时、员工满意度、重复咨询量
\end{itemize}
\end{infobox}


\section{总体架构（示意）}
\begin{figure}[H]
\centering
\begin{tikzpicture}[
  node distance=12mm and 12mm,
  box/.style={draw=Primary!40, rounded corners=2pt, fill=BoxBg, align=left, inner sep=6pt},
  svc/.style={draw=Accent!60, rounded corners=2pt, fill=Accent!6, align=left, inner sep=6pt},
  arr/.style={-Latex, draw=Primary!55, line width=0.7pt}
]
  \node[box] (user) {用户/员工\\(APP/小程序/坐席)};
  \node[box, right=of user] (entry) {业务入口\\(客服工作台/IM/门户)};
  \node[svc, right=of entry] (orchestrator) {编排层\\(提示词/路由/工具调用)};
  \node[svc, below=of orchestrator] (llm) {大模型推理服务\\(对话/理解/生成)};
  \node[box, right=of orchestrator] (systems) {客户系统\\(工单/CRM/订单/知识库)};
  \node[box, below=of systems] (kb) {知识库与数据源\\(文档/FAQ/日志)};

  \draw[arr] (user) -- (entry);
  \draw[arr] (entry) -- (orchestrator);
  \draw[arr] (orchestrator) -- node[above,sloped]{调用} (systems);
  \draw[arr] (orchestrator) -- node[right]{推理/生成} (llm);
  \draw[arr] (kb) -- node[below,sloped]{检索增强} (orchestrator);
  \draw[arr] (systems) -- node[right]{事件/结果} (orchestrator);
\end{tikzpicture}
\caption{大模型能力与客户系统集成总体架构示意（模板）}
\label{fig:arch}
\end{figure}

\section{部署、安全与运维要点（模板）}
\subsection{部署模式与数据范围}
\begin{itemize}[leftmargin=2.2em]
  \item 部署模式：公有云API/专有化（按实际选择，示例为「混合」）
  \item 主要数据源：知识库文档、工单文本、客服对话记录（脱敏后）
  \item 安全与合规：数据分级分类、脱敏与审计、权限隔离、日志留存
\end{itemize}

\subsection{运维与运营指标}
\begin{itemize}[leftmargin=2.2em]
  \item SLA：待确认（示例：7x24，关键链路可用性 99.9\%）
  \item 监控项：调用量/时延、错误率、命中率/引用率、越权/敏感词拦截
  \item 成本模型：按调用量/并发/部署资源计费（以合同为准）
\end{itemize}

\section{验收建议（模板）}
\begin{infobox}{建议的验收维度}
\begin{itemize}[leftmargin=2.2em]
  \item 功能验收：场景覆盖、工具调用链路、权限与审计、引用溯源（RAG）。
  \item 指标验收：自助解决率、时延、人工转接率、准确率（按业务定义）。
  \item 合规验收：敏感信息处理、越权防护、日志留存、数据出境（如适用）。
\end{itemize}
\end{infobox}

\end{document}
